\documentclass{covermaker}
\doctitle{Project Plan}
\docsubtitle{Scope, milestones, risks and way of working for Cover\_Maker v2}
\docid{CM-PP}
\docstatus{Released}
\input{glossary}

\begin{document}
\cmtitlepage

\begin{revisionhistory}
  0.1 & 2026-11-03 & Initial plan approved at the start of the v2 rewrite (M0). \\
  1.0 & 2026-11-05 & Version 1.0: all milestones M0--M12 delivered; scope, backlog and status updated. \\
  1.1 & 2026-11-10 & Re-plan after a competitive review: visual identity (M13) before the 1.0
  release, practice tools (M14), cover studio (M15), rehearsal and setlists (M16), launch plan
  and new risks R-09 to R-13. \\
  1.2 & 2026-11-11 & The pending release becomes 0.4.0; version 1.0.0 is kept for when
  M14, M15 and M16 are done. \\
\end{revisionhistory}

\tableofcontents

\chapter{Introduction}

\section{Purpose}
This document defines the scope, approach, milestones, risks and way of working
for version 2 of \CM{}. It is updated at the end of every milestone.

\section{Background}
Version 0.1 of \CM{} was a command-line tool that read a playlist from a music
streaming service, looked up each track on a video platform, downloaded the
audio and split it into two \glspl{stem} (vocals and instrumental). A review of
that version found functional defects (identity based on file names, wrong
matches, no tests) and, more importantly, significant legal risk in acquiring
audio from streaming platforms. Version 2 is a complete rewrite with a new
product definition; the old history was archived privately. Automatic lookup of
music on streaming services was dropped; downloading the audio of video links
pasted by the user was kept as an explicit, documented product decision
(\autoref{ch:legal}).

In November 2026 a competitive review found StemKit, an MIT-licensed desktop
application published in August 2026 that also splits songs from pasted video
links into stems. No code was shared in either direction, but its interface
looked very similar to \CM{}'s at the time. The plan was therefore revised
(\autoref{sec:positioning}): a visual identity of its own before the next
release, and a product that goes beyond separation to cover the whole path from
a song to a recorded cover.

\section{Product vision}
\cmnote{\CM{}, \emph{the cover musician's studio}, is a Windows desktop
application that separates local audio files or the audio of pasted video links
into the stems the user chooses (vocals, guitar, drums, bass, piano, other) or
renders the ``track without X''. Musicians then practise it in their key and
tempo, record their cover in audio and video, and rehearse their setlists.
Everything runs locally.}

\section{Positioning}\label{sec:positioning}
Splitting a song into stems is no longer a differentiator on its own. \CM{}
competes as the tool a cover musician uses from start to finish:
\begin{description}
  \item[Separate.] Choose the stems or the track without your instrument
        (version 0.4.0).
  \item[Practise.] Tuning, key and tempo analysis, pitch and tempo change
        without artefacts, bar-snapped loops, count-in and a speed trainer (M14).
  \item[Record.] Takes and webcam video over the backing track, in sync,
        exported for video platforms; the separated stems serve as guide
        tracks and an ``only my takes'' export keeps the original master out of
        anything published (M15).
  \item[Perform.] Per-song settings, setlists, a rehearsal/live mode and MIDI
        foot-switch control (M16).
\end{description}
Its other strengths remain: legal by design, local processing without
telemetry, verified model downloads with visible licences, \gls{daw}
integration, English and Spanish, and complete documentation.

\section{Target users}
Musicians who practise or record covers, and audio producers working on a
mid-range Windows PC or laptop, usually with a dedicated GPU and often with a
\gls{daw} open at the same time.

\chapter{Scope}

\section{In scope for version 0.4.0}
\begin{itemize}
  \item Selection of stems to extract or to remove (``track without X''), with
        the option to keep or drop the \gls{residual}.
  \item Local files by drag and drop, and watched folders processed incrementally.
  \item Video and playlist links pasted by the user, with audio-only download
        through yt-dlp and a rights confirmation per job.
  \item Preview mixer with synced waveforms, mute, solo and volume per stem.
  \item Job queue with per-song progress, pause, cancel and retry.
  \item Library with re-sync: only new or changed inputs are processed again.
  \item Import of playlist files (streaming data export JSON, CSV, JSPF, M3U)
        matched against the local library, with manual match correction.
  \item Creative Commons catalogues with licence display and attribution.
  \item Model catalogue with licence class and provenance; no default model.
  \item Hardware detection, presets per hardware tier and a low-impact mode
        for \gls{daw} coexistence.
  \item English and Spanish user interface; installer with auto-update.
  \item \gls{daw} integration: a REAPER script that inserts each stem as a new
        track and a VST3 plugin that previews the stems and lets the user drag
        them into tracks (M11, brought forward before 1.0).
\end{itemize}

\section{Planned before version 1.0}
Version 0.4.0 ships after M13, the new visual identity, which also brings
per-stem pan, keyboard shortcuts and the interface animations. Then follow
practice tools (M14), the cover studio for recording audio and video (M15) and
rehearsal with setlists (M16), each in its own minor release;
\autoref{ch:milestones} lists their goals. Version 1.0.0 is released only when
all three are done.

\section{Backlog}
\gls{daw} session export (REAPER and DAWproject with tempo map, markers and
key), BWF metadata and LUFS-normalised exports, batch and headless processing,
audio-to-MIDI per stem, chord detection, lead/backing vocal separation, drum
sub-stems, de-reverb, a full mini \gls{daw} with ASIO recording, more interface
languages, macOS, metadata and cover-art tagging, ONNX/DirectML acceleration
for AMD and Intel GPUs, and cloud separation backends. Licences are checked
before any of these is scheduled: components under GPL, AGPL or
non-commercial terms are avoided.

\section{Explicitly excluded}
Searching, scraping or matching music on streaming services, use of their
\glspl{api} to locate music, and decryption of protected (DRM) streams (see
\autoref{ch:legal}).

\chapter{Legal by design}\label{ch:legal}
A risk analysis (not legal advice) carried out before the rewrite concluded that
stem separation of lawfully obtained audio carries very low risk, while
acquiring audio from streaming platforms does not: in Spain the private-copy
exception requires a lawful source and cannot override technological protection
measures, German courts have treated a major video platform's stream protection
as an effective \gls{tpm}, and a large streaming service's developer terms
forbid using its metadata to obtain music elsewhere. The full analysis with
sources is kept in the Compliance chapter of the Architecture \& Design
Document.

On 2026-11-03 the product owner decided to keep the download of audio from
video links pasted by the user, using the unmodified yt-dlp project as a
dependency. The residual risk is accepted and tracked as R-08: if a takedown
notice is received, the project will follow the approach of the yt-dlp project
and keep a disclaimer, and can move the feature into a separate optional plugin.
The following constraints are binding for the whole project:

\input{generated/constraints}

\chapter{Target hardware}
The reference machine is a mid-range producer PC or laptop. Many PCs sold for
audio work have no dedicated GPU, so CPU operation is a first-class case.

\begin{table}[h]
\centering
\begin{tabularx}{\linewidth}{@{}lYYY@{}}
\toprule
 & \textbf{Minimum} & \textbf{Recommended} & \textbf{Ideal} \\
\midrule
CPU & 4 cores / 8 threads, AVX2 & 6--8 cores, 2022 or newer & 8+ cores \\
RAM & 16 GB & 32 GB & 32--64 GB \\
GPU & Optional; NVIDIA Turing+ with 4 GB & NVIDIA RTX, 8--12 GB & NVIDIA RTX, 12--16 GB \\
Disk & SSD, 10 GB free & NVMe, 20 GB+ free & NVMe \\
\bottomrule
\end{tabularx}
\caption{Hardware reference specification.}
\end{table}

Estimated separation times per hardware tier are recorded in the Architecture
\& Design Document and will be replaced by measurements from milestone M1. The
development machine (RTX 3050 Laptop, 4 GB \gls{vram}) sits on the minimum GPU
tier.

\chapter{Approach}

\section{Lifecycle}
The project is incremental and risk-driven. The two riskiest technical unknowns
(model performance on low-\gls{vram} hardware, and the desktop shell lifecycle
and first-run bootstrap) are tackled as spikes in M1 and M2 before the
production engine is built. Every milestone ends with something usable or
demonstrable. There is no fixed deadline; effort is estimated relatively
(S, M, L).

\section{Workflow}
\begin{enumerate}
  \item A GitHub milestone and its issues are created when work on it starts.
  \item Work happens on a branch named \texttt{mX/<slug>}, pushed frequently.
  \item A pull request is opened with green \gls{ci}; it is merged after the
        product owner's approval and tagged as a pre-release.
\end{enumerate}

\section{Definition of done}
Tests added and passing; lint and type checks clean; affected documents and the
changelog updated in the same pull request; a reproducible demo described in the
pull request; legal-by-design checklist ticked.

\section{Tooling}
\begin{tabularx}{\linewidth}{@{}lY@{}}
\toprule
\textbf{Area} & \textbf{Tools} \\
\midrule
Engine & Python 3.12, uv, ruff, pyright (strict), pytest \\
Desktop & Tauri 2, Rust (MSVC), SvelteKit + TypeScript, pnpm, eslint, prettier, vitest \\
Documentation & LaTeX (LuaLaTeX, latexmk), YAML sources rendered by \texttt{gen\_tex.py} \\
Delivery & GitHub Actions (Windows for code, Ubuntu for documents), GitHub Releases \\
\bottomrule
\end{tabularx}

\chapter{Milestones}\label{ch:milestones}
\input{generated/milestones}

\section{Launch plan}
The repository and the download are treated as part of the product. Before
1.0: a new README in English and Spanish with a short demonstration video and a
download button, a stable download link to the latest installer, a signed
installer, package-manager entries (winget, Scoop), a landing page on the
author's website and a bundled Creative Commons example song so that a first
separation takes minutes. The launch then runs in waves: the Spanish-speaking
community with 1.0, practice-focused short videos with M14, and the main launch
(technology forums, product directories, music communities and the press) with
M15. Promotional material uses only the author's own or CC-licensed music, and
downloading from video platforms is not used as the hook (risks R-08 and
R-13).

\chapter{Risk management}
Risks are reviewed at the end of each milestone.
\input{generated/risks}

\chapter{Documentation plan}
\begin{tabularx}{\linewidth}{@{}lYll@{}}
\toprule
\textbf{Document} & \textbf{Basis} & \textbf{Starts} & \textbf{Closes} \\
\midrule
Project Plan & Simplified PMBOK & M0 & every milestone \\
Software Requirements Specification & ISO/IEC/IEEE 29148 & M0 & M8 \\
Architecture \& Design Document & arc42, C4, ADRs & M0 & M9 \\
Test Plan \& Report & ISO/IEC/IEEE 29119-3 (light) & M3 & M10 \\
Installation Guide & -- & M2 & M9 \\
User Manual & -- & M3 & M10 \\
\bottomrule
\end{tabularx}

\chapter{Decision log}
\begin{tabularx}{\linewidth}{@{}lY@{}}
\toprule
\textbf{Topic} & \textbf{Decision (2026-11-03)} \\
\midrule
Distribution & Public open source, MIT licence, Windows only. \\
Repository & Same repository; v0.1 history archived privately and replaced. \\
Desktop stack & Tauri 2 + Svelte 5 shell, Python engine over local HTTP + WebSocket. \\
Audio sources & Local files, watched folders, playlist-file import, CC catalogues, user-pasted video links (yt-dlp). \\
Streaming services & No search, no API use, no scraping, no DRM decryption. \\
Video links (2026-11-03) & Kept in scope; risk R-08 accepted; rights confirmation and disclaimer. \\
Models & No default model; user chooses and accepts each licence. \\
Documentation & LaTeX in English, six documents, built in CI. \\
Interface & English and Spanish. \\
Competition (2026-11-09) & Position \CM{} as the cover musician's studio:
own visual identity, practice tools, cover recording and setlists. \\
Visual identity (2026-11-09) & ``Tape \& Console'', a warm hardware look, in M13. \\
Name (2026-11-09) & \CM{} with the descriptor ``the cover musician's studio'';
trademark search before marketing. \\
Release 1.0 (2026-11-09) & Published after M13, with the new identity. \\
Launch (2026-11-09) & In waves; own or CC-licensed music only in promotional material. \\
\bottomrule
\end{tabularx}

\printnoidxglossaries
\end{document}
